% ========================================================
% chapters/theory/genus_2_ddybe_platonic.tex
%
% Structural inscription of the genus-2 doubly-dynamical
% Yang-Baxter equation. Separates the generic Omega
% weight-zero check at 10^{-12} from the full engine tolerance 1e-4
% and disambiguates the two statements called
% "Fay trisecant" (Szego three-term identity universal at all
% genera vs Riemann theta-nulls quartic identity on the theta
% divisor).
%
% ========================================================

\providecommand{\cA}{\mathcal{A}}
\providecommand{\cC}{\mathcal{C}}
\providecommand{\cL}{\mathcal{L}}
\providecommand{\cM}{\mathcal{M}}
\providecommand{\CC}{\mathbb{C}}
\providecommand{\ZZ}{\mathbb{Z}}
\providecommand{\RR}{\mathbb{R}}
\providecommand{\HHH}{\mathbb{H}}
\providecommand{\fh}{\mathfrak{h}}
\providecommand{\fsl}{\mathfrak{sl}}
\providecommand{\Conf}{\mathrm{Conf}}
\providecommand{\FM}{\overline{\mathrm{FM}}}
\providecommand{\End}{\mathrm{End}}
\providecommand{\ClaimStatusProvedHere}{\textsf{ProvedHere}}
\providecommand{\ClaimStatusProvedElsewhere}{\textsf{ProvedElsewhere}}
\providecommand{\ClaimStatusConjectured}{\textsf{Conjectured}}
\providecommand{\ClaimStatusNumericalEvidence}{\textsf{NumericalEvidence}}

\chapter[Genus-$2$ doubly-dynamical Yang-Baxter]
{Genus-$2$ doubly-dynamical Yang-Baxter equation}
\label{ch:genus-2-ddybe-platonic}

\section*{Preface to this chapter}

The genus-$2$ construction contains two logically separate facts.
The degree-$2$ ordered chiral homology on a Torelli-marked genus-$2$
surface is computed in Proposition~\ref{prop:g2-degree2} and gives
Euler characteristic $\chi=-12$. The generic-$\Omega$
doubly-dynamical Yang--Baxter equation (DDYBE) is weaker: the
$10^{-12}$ agreement belongs to the single-entry weight-zero check,
while the full three-term identity is evaluated by the engine
(\texttt{face\_model\_ddybe\_engine.py},
\texttt{verify\_face\_ddybe\_g2}) only to tolerance
$\mathrm{relative}<10^{-4}$. The chapter keeps these facts on
separate verification surfaces.

Four statements form the genus-$2$ package:
\begin{enumerate}[label=\textup{(\arabic*)}]
\item \textbf{Existence of the flat KZB connection at $g=2$}
 (\cref{thm:genus-2-kzb-connection-platonic}): the
 Calaque--Enriquez--Etingof universal construction specialises
 to $\mathfrak{sl}_2$ evaluation modules as a flat connection
 on $\overline{\cM}_{2,n}\times\HHH_2$ with three independent
 modular directions $\Omega_{11}, \Omega_{12}, \Omega_{22}$.
 \emph{Verification: CEE09.}
\item \textbf{Fay trisecant identity, correct name}
 (\cref{thm:fay-trisecant-genus-2-specific}): Fay's 1973
 Corollary~2.5 is the normalised Szeg\H{o}/prime-form
 identity on $C_3(X)$ for all $g\ge 0$; in the torus
 normalisation it contains the load-bearing
 $\vartheta_1'(0,\tau)^2$ factor. A distinct four-factor
 theta-null product identity holds only on the theta divisor;
 the two statements are not interchangeable.
 \emph{Verification: Fay 1973.}
\item \textbf{Face-model DDYBE: exact degeneration plus finite-window evidence}
 (\cref{thm:g2-face-model-bypass-scope-restricted,prop:g2-generic-ddybe-finite-window-evidence}):
 the diagonal/separating face-model degeneration is a theorem,
 because the genus-$2$ theta series factorises and the equation
 reduces to a genus-$1$ Felder DYBE. Generic off-diagonal
 $\Omega$ checks are numerical evidence only: relative residual
 $<10^{-4}$ at finitely many parameter points ($4$ generic-$\Omega$
 compute tests, one second-$\Omega$ spot check, plus a tier-(T3)
 diagonal-$\Omega$ consistency check).
 \emph{Verification: exact on the diagonal/separating locus;
 numerical finite-window evidence at generic \(\Omega\) for the
 conjectural full DDYBE.}
\item \textbf{Euler characteristic at degree~$2$}
 (\cref{cor:g2-chi-minus-12}): $\chi=-12$ follows from the
 rank-$4$ KZB local system on $\Sigma_2\setminus\{0\}$ via
 $\chi=r(2-2g-s)$. This is a topological degree-counting
 invariant, not a doubly-dynamical integrable-system
 invariant.
 \emph{Status: ProvedHere.}
\end{enumerate}

The conjecture \texttt{conj:g2-ddybe} retains its
\texttt{\ClaimStatusConjectured} status.

\section{Setup: Torelli-marked genus-$2$ surface}
\label{sec:g2-ddybe-setup}

Let $\Sigma_2$ be a smooth projective curve of genus~$2$, and
fix a \emph{Torelli marking}: a symplectic basis
$(A_1, A_2, B_1, B_2)$ of $H_1(\Sigma_2,\ZZ)$ with
$A_i\cdot B_j = \delta_{ij}$, $A_i\cdot A_j = B_i\cdot B_j = 0$.
Let $\omega_1, \omega_2 \in H^0(\Sigma_2, K_{\Sigma_2})$ be the
uniquely determined normalised holomorphic differentials
with $\oint_{A_i}\omega_j = \delta_{ij}$. The \emph{Siegel
period matrix} is
\begin{equation}\label{eq:g2p-siegel}
 \Omega \;=\;
 \begin{pmatrix}
 \Omega_{11} & \Omega_{12} \\
 \Omega_{12} & \Omega_{22}
 \end{pmatrix},
 \qquad
 \Omega_{ij} \;=\; \oint_{B_i}\omega_j,
\end{equation}
which lies in the Siegel upper half-space
\begin{equation}\label{eq:g2p-H2}
 \HHH_2 \;=\;
 \bigl\{ Z \in M_2(\CC) \;:\; Z = Z^t,\;
 \mathrm{Im}(Z) > 0 \bigr\}.
\end{equation}
The symmetry $\Omega_{12}=\Omega_{21}$ is the Torelli theorem
at $g=2$ (a consequence of Riemann bilinear relations); the
positive-definiteness of $\mathrm{Im}(\Omega)$ is the Riemann
bilinear inequality.

The prime form and Szeg\H{o} kernel on $\Sigma_2$ are defined in
Equations~(\ref{eq:prime-form-g2}) and
(\ref{eq:szego-g2}) of the main chapter
\texttt{higher\_genus\_modular\_koszul.tex}; we do not reproduce
them here.

\begin{remark}[separating versus non-separating
degenerations]
\label{rem:g2-ddybe-check}
The boundary divisor $\partial\overline{\cM}_2$ has two
irreducible components.
\begin{enumerate}[label=\textup{(\alph*)}]
\item \emph{Non-separating (irreducible)} $\delta_{\mathrm{irr}}$:
 a single non-separating vanishing cycle on $\Sigma_2$
 collapses, producing an irreducible nodal curve of
 arithmetic genus~$2$. Geometrically, one handle of
 $\Sigma_2$ pinches. Algebraically:
 $\Omega_{22}\to i\infty$, $\Omega_{12}\to 0$,
 $\Omega_{11}$ fixed (up to $SL_2(\ZZ)$ action).
\item \emph{Separating} $\delta_{\mathrm{sep}}$: a separating
 cycle collapses, producing a nodal curve whose normalisation
 is a disjoint union $E_{\tau_1}\sqcup E_{\tau_2}$ of two
 elliptic curves. Algebraically: $\Omega_{12}\to 0$
 with $\Omega_{11}=\tau_1$, $\Omega_{22}=\tau_2$ finite.
\end{enumerate}
The separating regime carries no genuinely
genus-$2$ DDYBE content (a degeneration-dependent identity is
not a genus-$2$ invariant), because the genus-$2$ theta function
factorises as
$\Theta_{\Sigma_2}(z_1, z_2 \mid
\mathrm{diag}(\tau_1,\tau_2))
= \theta_1(z_1\mid\tau_1)\cdot\theta_1(z_2\mid\tau_2)$
(up to a sign from the odd characteristic), and the DDYBE
reduces on the nose to two independent copies of the
Felder genus-$1$ DYBE. The genuinely genus-$2$ content of
\texttt{conj:g2-ddybe} lives in (i) the generic $\Omega$
interior of $\HHH_2$ and (ii) the non-separating-boundary
limit, where the off-diagonal period $\Omega_{12}$ survives
the limit $\Omega_{22}\to i\infty$.
\end{remark}

% ========================================================
\section{Flat KZB connection at $g=2$}
\label{sec:g2-kzb-flat}

\begin{theorem}[Flat KZB connection on $\overline{\cM}_{2,n}\times\HHH_2$;
\ClaimStatusProvedElsewhere]
\label{thm:genus-2-kzb-connection-platonic}
\index{KZB connection!genus $2$}
\index{Calaque-Enriquez-Etingof!universal KZB}
Let $n\ge 1$ and let $\fsl_2$-evaluation modules $V_1,\ldots,V_n$
be attached to $n$ ordered marked points on $\Sigma_2$. The
Calaque--Enriquez--Etingof universal
KZB connection~\cite{CEE09} specialises to a flat connection
\begin{equation}\label{eq:g2p-kzb-flat}
 \nabla^{\mathrm{KZB}}_{(g=2,n)} \;=\;
 d_{\mathrm{spatial}} \;+\;
 d_{\mathrm{modular}}
\end{equation}
on the trivial local system
$\underline{V_1\otimes\cdots\otimes V_n}$
over an open dense locus of
$\Conf_n^{\mathrm{ord}}(\Sigma_2)\times\HHH_2$, with
\begin{align}
 d_{\mathrm{spatial}}
 &\;=\; d \;-\; \hbar \sum_{i<j} S_2(z_i,z_j)\,\Omega_{ij},
 \label{eq:g2p-kzb-spatial} \\
 d_{\mathrm{modular}}
 &\;=\; \sum_{\alpha\le \beta}
 \Bigl(\partial_{\Omega_{\alpha\beta}}
 -\tfrac{1}{2\pi i}\hbar\sum_{i<j}
 \partial_{z_i}\partial_{z_j}\log E(z_i,z_j)\,
 \omega_\alpha(z_i)\omega_\beta(z_j)\,\Omega_{ij}\Bigr)
 d\Omega_{\alpha\beta},
 \label{eq:g2p-kzb-modular}
\end{align}
where $\hbar = 1/(k+h^\vee) = 1/(k+2)$, $S_2$ is the Szeg\H{o}
kernel, $E$ the prime form, and $\Omega_{ij}$ is the
$\fsl_2$-Casimir acting in positions $i,j$.
The modular part has three independent components
$(\alpha,\beta)\in\{(1,1),(1,2),(2,2)\}$: the Siegel upper
half-space has complex dimension~$3$. The connection is flat:
\begin{equation}\label{eq:g2p-flatness}
 \bigl(\nabla^{\mathrm{KZB}}_{(g=2,n)}\bigr)^2 \;=\; 0.
\end{equation}
\end{theorem}

\begin{proof}[Proof sketch and attribution]
The universal flat connection on
$\Conf_n(\Sigma_g)\times\mathrm{Teich}_g$ is constructed by
Calaque-Enriquez-Etingof~\cite{CEE09} as the image of a formal
connection with values in the Chen-integral completion of the
derivation Lie algebra of the genus-$g$ surface group. Its
flatness is proved by a direct calculation on the universal
cover and transported to the configuration space by
pull-back (\cite[Theorems 9.3, 9.11]{CEE09}). Specialising to
$n=2$ and to $\fsl_2$-evaluation
modules $V^{\otimes n}$ yields the stated system; the
Szeg\H{o}-kernel presentation of the spatial part uses
Bergman-kernel technology from
Fay~\cite[Ch.~1]{Fay73}. The three-direction modular part is
the special case of the CEE universal construction at $g=2$,
where $\dim_\CC\HHH_2 = 3$. The heat-equation identity
$\partial_{\Omega_{\alpha\beta}}\theta[\delta]
= \tfrac{1}{2\pi i}\partial_{z_\alpha}\partial_{z_\beta}
\theta[\delta]$ enforces compatibility of spatial and modular
components; its specialisation to the $\fsl_2$ KZB system is
direct.
\end{proof}

\begin{remark}[Genus-$1$ limit]
\label{rem:g2p-g1-limit}
At the non-separating boundary $\Omega_{22}\to i\infty$,
$\Omega_{12}\to 0$, the connection
\eqref{eq:g2p-kzb-flat} reduces to the
Bernard~\cite{Bernard88} genus-$1$ KZB connection on
$\Conf_n(E_{\Omega_{11}})\times\HHH_1$ with $\tau=\Omega_{11}$.
At the separating boundary $\Omega_{12}\to 0$ with
$\Omega_{11},\Omega_{22}$ finite, the connection splits
(via $\omega_1\to dz$ on $E_{\Omega_{11}}$ and
$\omega_2\to dz$ on $E_{\Omega_{22}}$) into a product of two
independent genus-$1$ KZB connections; this is the
factorised-limit regime, which carries no genuinely $g=2$
content.
\end{remark}

% ========================================================
\section{Fay trisecant: the correct identification}
\label{sec:g2-fay}

\begin{theorem}[Fay trisecant, normalised Szeg\H{o}/prime-form form;
\ClaimStatusProvedElsewhere]
\label{thm:fay-trisecant-genus-2-specific}
\index{Fay trisecant identity!normalised Szeg\H{o} form}
Let $X$ be a smooth projective curve of any genus $g\ge 0$,
equipped with a choice of odd non-singular theta
characteristic $\delta$ where the Szeg\H{o} kernel is defined.
Fay's Corollary~2.5 is a theta-addition identity for the
normalised prime-form/Szeg\H{o} kernels on $C_3(X)$; after
pullback to a three-point Fulton--MacPherson boundary and
contraction with the Arnold--Kohno/IHX tensor relation, it gives
the Stokes cancellation used by the bar complex. It is not the
bare scalar equality
$S_{12}S_{23}+S_{23}S_{31}+S_{31}S_{12}=0$ for an arbitrarily
normalised scalar kernel.

In the genus-$1$ normalisation with periods $(1,\tau)$,
$\vartheta_1=\vartheta_1(-\mid\tau)$, and Weierstrass
$\wp=\wp(-,\tau)$, the torus specialisation is
\begin{equation}\label{eq:fay-three-term}
 \wp(a,\tau)-\wp(b,\tau)
 =
 -\,\frac{
 \vartheta_1(a+b,\tau)\,\vartheta_1(a-b,\tau)}
 {\vartheta_1(a,\tau)^2\,\vartheta_1(b,\tau)^2}
 \,\vartheta_1'(0,\tau)^2 .
\end{equation}
At $g=0$ the same identity degenerates to the rational
partial-fraction relation. At $g=2$ the corresponding statement is
the full prime-form/theta-characteristic identity on
$\operatorname{Jac}(X)$, with the same normalisation discipline:
theta-prime and prime-form factors are part of the identity.
\end{theorem}

\begin{proof}[Proof and attribution]
This is Corollary~2.5 of Fay~\cite{Fay73}, proved uniformly
for all $g\ge 0$ via the addition formula for theta
functions on $\mathrm{Jac}(X)$. In genus~$1$ the displayed
formula is the standard theta-prime-normalised torus form, also
recorded in Mumford's \emph{Tata Lectures on Theta I}, Ch.~I
\S5. The Szeg\H{o}-kernel simple-pole structure is independent
of the theta-function construction (Bergman 1958, Fay 1973
Ch.~1), providing a verification pathway disjoint from the
theta-addition route.
\end{proof}

\begin{remark}[Name disambiguation: Szeg\H{o}/prime-form vs
theta-nulls quartic]
\label{rem:fay-name-disambiguation}
Two distinct identities are called ``Fay trisecant'' in the
literature; the bar-complex argument uses the normalised
Szeg\H{o}/prime-form identity whose torus form is
\eqref{eq:fay-three-term}.
\begin{enumerate}[label=\textup{(\alph*)}]
\item \emph{Normalised Szeg\H{o}/prime-form form}
 (\eqref{eq:fay-three-term}, Fay 1973 Cor.~2.5): universal
 on $C_3(X)$ at all $g\ge 0$ after the prime-form and
 theta-characteristic normalisations; no divisor condition.
\item \emph{Four-theta-nulls quartic identity}
 (Fay 1973 Cor.~4.8 and variants, ``trisecant formula'' in
 the Schottky-locus literature): a product of four theta
 values on the Jacobian satisfies a quartic relation
 whose vanishing characterises the theta divisor (hence
 the name ``trisecant''). This identity HOLDS ONLY ON THE
 THETA DIVISOR, not identically on $\HHH_g$; it is a
 divisor-condition identity, not a kernel identity.
\end{enumerate}
The bar-differential identity $d_{\mathrm{bar}}^2=0$,
the doubly-dynamical YBE, and all Stokes-regularity
arguments use form~(a), not form~(b). Any reference to
``Fay trisecant'' in the prose of Vol~I refers to
form~(a) unless explicitly stated otherwise.
\end{remark}

\begin{remark}[Why (a), not (b), appears in the bar complex]
\label{rem:why-a-not-b}
The bar differential $d_{\mathrm{bar}}$ on
$\Omega^*_{\log}(\FM_n^{\mathrm{ord}}(\Sigma_2))\otimes
(s^{-1}\bar\cA)^{\otimes n}$ receives a three-propagator
Stokes term at each 3-boundary
$D_{\{i,j,k\}}\subset\FM_n$. The vanishing of this Stokes
term is the section-level consequence of Fay's normalised
prime-form/Szeg\H{o} identity, together with the
Arnold--Kohno/IHX tensor contraction. The unnormalised scalar
product $S_{12}S_{23}+S_{23}S_{31}+S_{31}S_{12}$ is not used as
a standalone equality. The theta-nulls quartic identity would
govern a four-point interaction only on the Schottky locus, which
is a proper codimension subvariety of $\HHH_2$; the bar complex
is defined over the full $\HHH_2$ interior, so only identity~(a)
suffices.
\end{remark}

% ========================================================
\section{Face-model DDYBE: exact degeneration and numerical evidence}
\label{sec:g2-face-model-bypass}

The Baxter--Felder face-model construction replaces the
vertex $R$-matrix $R(z,\lambda)\in\End(V\otimes V)$ by the
four-index Boltzmann weights
$w(a,b,c,d\mid z, \lambda)$ built from theta functions
directly. At genus~$2$, the vertex-IRF transform is not
established in the literature, so the face model is
constructed on its own from genus-$2$ theta-Boltzmann weights;
this is the operational meaning of ``bypass''.

\begin{definition}[Genus-$2$ face-model Boltzmann weights]
\label{def:g2-face-boltzmann}
\index{Boltzmann weights!genus-$2$ face model}
Fix $\eta=\hbar=1/(k+h^\vee)$. At the odd theta characteristic
$\delta=({}^{1/2,\,0}_{1/2,\,0})$, write
$\Theta_\delta(z\mid\Omega)$ for the corresponding genus-$2$
Riemann theta function. The \emph{face-model Boltzmann
weights} at genus~$2$ are the four entries
\begin{equation}\label{eq:g2-boltzmann}
\begin{split}
 w_{\text{prop}}(a\mid z,\Omega)
 &=
 \frac{\Theta_\delta(z\,e_1\mid\Omega)\,\Theta_\delta((a+1)\eta\,e_1\mid\Omega)}
 {\Theta_\delta(\eta\,e_1\mid\Omega)\,\Theta_\delta((a\eta+z)e_1\mid\Omega)}, \\
 w_{\text{refl}}(a\mid z,\Omega)
 &=
 \frac{\Theta_\delta(a\eta\,e_1\mid\Omega)\,\Theta_\delta(z+\eta\,e_1\mid\Omega)}
 {\Theta_\delta(\eta\,e_1\mid\Omega)\,\Theta_\delta((a\eta+z)e_1\mid\Omega)},
\end{split}
\end{equation}
with $e_1=(1,0)^t$ the first coordinate direction, and two
further configurations related by the symmetries
$a\mapsto -a$, $z\mapsto -z$. The composite face-model
$R$-matrix is the $4\times 4$ matrix
\begin{equation}\label{eq:g2-face-R}
 R^{\mathrm{face}}(z,\boldsymbol\lambda,\Omega)
 \;\in\; \End(V\otimes V),
 \qquad V=\CC^2,
\end{equation}
whose entries are linear combinations of
the four Boltzmann weights, with
$\boldsymbol\lambda=(\lambda_1,\lambda_2)\in\fh^*\oplus\fh^*$.
\end{definition}

\begin{theorem}[Exact diagonal/separating degeneration to Felder DYBE;
\ClaimStatusProvedHere]
\label{thm:g2-face-model-bypass-scope-restricted}
\index{doubly-dynamical Yang-Baxter equation!face-model,
exact degeneration}
The face-model $R$-matrix $R^{\mathrm{face}}$ of
Definition~\ref{def:g2-face-boltzmann}
satisfies the exact diagonal/separating degeneration statements
below. These statements are theorem-level because they consume only
theta-series factorisation and the genus-$1$ Felder
dynamical Yang--Baxter theorem; they do not prove the full
generic-$\Omega$ system \eqref{eq:ddybe}--\eqref{eq:ddybe-coupling}.
\begin{enumerate}[label=\textup{(\roman*)}]
\item \textbf{Diagonal $\Omega$ regime (exact):}
 at $\Omega=\mathrm{diag}(\tau_1,\tau_2)$ with
 $\Omega_{12}=0$ and $\tau_1,\tau_2\in\HHH_1$, the
 Boltzmann weights
 \eqref{eq:g2-boltzmann} use $\Theta_\delta$ evaluated along
 the $e_1$-direction only; at diagonal $\Omega$ the series
 factorises as
 $\Theta_\delta(x\,e_1\mid\mathrm{diag}(\tau_1,\tau_2))
 = \epsilon_\delta\,\theta_1(x\mid\tau_1)\,
 \theta_3(0\mid\tau_2)$
 with $\epsilon_\delta=\pm 1$, and the constant factor
 $\epsilon_\delta\,\theta_3(0\mid\tau_2)$ cancels in every
 Boltzmann-weight ratio~\eqref{eq:g2-boltzmann}, so
 \begin{equation}\label{eq:g2p-diagonal-exact}
 R^{\mathrm{face}}\bigl(z,(\lambda_1,\lambda_2),
 \mathrm{diag}(\tau_1,\tau_2)\bigr)
 \;=\;
 R^{\mathrm{Felder}}(z,\lambda_1,\tau_1),
 \end{equation}
 independent of $\lambda_2$ and $\tau_2$: the face-model
 $R$-matrix reduces to a single genus-$1$ Felder $R$-matrix,
 and the DDYBE~\eqref{eq:ddybe} reduces on the nose to a
 single instance of the Felder~\cite{Felder94} dynamical
 Yang-Baxter equation, which is a theorem
 (Felder 1994, Etingof-Varchenko 1998). In particular the
 diagonal specialisation does not witness nontrivial doubly-
 dynamical content; it is a consistency check against the
 genus-$1$ theorem.
\item \textbf{Separating-degeneration regime
 (factorised-limit, no genuinely $g=2$ content):} at the
 separating boundary $\Omega=\mathrm{diag}(\tau_1,\tau_2)$,
 the face $R$-matrix reduces by~(i) to a single genus-$1$
 Felder $R$-matrix and the DDYBE reduces to a single Felder
 DYBE; this regime carries no \emph{genuinely genus-$2$}
 DDYBE content.
\end{enumerate}
\end{theorem}

\begin{proof}
Substitute $\Omega=\mathrm{diag}(\tau_1,\tau_2)$ into the
genus-$2$ Riemann theta series
$\Theta_\delta(z\mid\Omega)=\sum_{n\in\ZZ^2+\delta_1}
e^{\pi i n^t\Omega n + 2\pi i n^t(z+\delta_2)}$. The sum
over $n=(n_1,n_2)$ separates into two independent sums in
$n_1,n_2$; with the odd characteristic
$\delta=({}^{1/2,0}_{1/2,0})$, the $n_1$-sum produces
$\theta_1(z_1\mid\tau_1)$ (odd) and the $n_2$-sum
produces $\theta_3(z_2\mid\tau_2)$ (even) up to an overall
$\epsilon_\delta$ sign fixed by $\delta_2$. Each Boltzmann
weight in~\eqref{eq:g2-boltzmann} is a ratio of theta values
at the SAME $z$-direction ($z e_1 + \text{shifts}$), so the
$\theta_3(\cdot\mid\tau_2)$ factor appears identically in
numerator and denominator and cancels. The cancellation is
\emph{algebraic}, taking place at the level of the theta series
BEFORE any lattice-truncation is applied; the numerical
regression gate
\texttt{verify\_g2\_to\_g1\_degeneration}
(\texttt{compute/lib/face\_model\_ddybe\_engine.py}, tested at
\texttt{test\_degeneration\_to\_g1} and
\texttt{test\_degeneration\_different\_params}) records
element-wise equality
$R^{\mathrm{face}}_{g=2}(z,\lambda_1,\eta,\mathrm{diag}(\tau_1,\tau_2))
= R^{\mathrm{face}}_{g=1}(z,\lambda_1,\eta,\tau_1)$
to relative residual $<10^{-6}$ at $N=10$, corroborating the
algebraic identity. The resulting
weights are exactly the genus-$1$ Felder Boltzmann weights
in the variable $z$ with parameter $\tau_1$, and the
composite $R$-matrix becomes
$R^{\mathrm{Felder}}(z,\lambda_1,\tau_1)$, independent of
$(\lambda_2,\tau_2)$. The DDYBE~\eqref{eq:ddybe} is now a
single instance of Felder's DYBE (Felder 1994 Theorem 3.1,
Etingof-Varchenko 1998 Theorem 1.1), which is a theorem.

The separating-degeneration statement is the same factorised-limit
calculation read as a boundary statement in
$\overline{\cM}_2$: the $\Omega_{12}\to0$ boundary has two elliptic
components, and the $e_1$-direction face weight sees only the
$\tau_1$ component after the common $\theta_3(0|\tau_2)$ factor is
cancelled. Thus the boundary equation is a single genus-$1$ Felder
DYBE and carries no genuinely genus-$2$ content.
\end{proof}

\begin{proposition}[Generic-\texorpdfstring{$\Omega$}{Omega} DDYBE
finite-window residuals; \ClaimStatusNumericalEvidence]
\label{prop:g2-generic-ddybe-finite-window-evidence}
\index{doubly-dynamical Yang-Baxter equation!generic genus-$2$ numerical evidence}
For generic $\Omega\in\HHH_2$ with $\Omega_{12}\ne 0$, the
face-model DDYBE residual is \(<10^{-4}\) in the finite windows
listed below. This is numerical evidence for
Conjecture~\textup{\ref{conj:g2-ddybe}}, not a proof of the
full generic-$\Omega$ DDYBE and not a uniform estimate on
\(\HHH_2\).

The engine \texttt{compute/lib/face\_model\_ddybe\_engine.py}
function \texttt{verify\_face\_ddybe\_g2} asserts
\texttt{result['passed']} under the threshold
\texttt{relative < 1e-4}, across four distinct
$(z,w,\boldsymbol\lambda,\eta)$ parameter tuples sharing the
single $\Omega$ value
$\Omega_0 = \begin{pmatrix}1.1\mathrm{i} & 0.15+0.05\mathrm{i}\\
0.15+0.05\mathrm{i} & 1.3\mathrm{i}\end{pmatrix}$
(two distinct $\boldsymbol\lambda$, two distinct $(z,w)$ spectral
pairs, $\eta=0.25, 0.2, 0.4$: levels $k=2,3$ and strong-coupling
$\eta=0.4$), plus a second-$\Omega$ spot check at
$\Omega_1 = \begin{pmatrix}2\mathrm{i} & 0.3+0.1\mathrm{i}\\
0.3+0.1\mathrm{i} & 1.5\mathrm{i}\end{pmatrix}$.
The sweep covers a positive-codimension subset of the
three-complex-dimensional moduli $\HHH_2$ and provides
pointwise, not uniform, evidence; the bulk of the sampling
concentrates at the single point $\Omega_0$, with a single
second-$\Omega$ witness. A sixth test at
$\Omega=\mathrm{diag}(1\mathrm{i},1.3\mathrm{i})$,
$\boldsymbol\lambda=(\lambda_1,0)$ is included as a diagonal-$\Omega$
sanity check at tier~(T3); it carries no genuinely $g=2$
content by Theorem~\ref{thm:g2-face-model-bypass-scope-restricted}.
\end{proposition}

\begin{proof}[Numerical verification protocol]
The 4x4 matrices $R^{\mathrm{face}}(z,\boldsymbol\lambda,\Omega)$
are evaluated at $N=8$ lattice-truncation terms for the
genus-$2$ Riemann theta, following the standard exponential
decay of the theta series with bound
$\varepsilon(N) \lesssim e^{-\pi N^2 \min(\mathrm{Im}(\Omega))/2}$
governing truncation error. The three matrix products LHS and
RHS of~\eqref{eq:ddybe} are computed and the relative residual
$\max|\mathrm{LHS}-\mathrm{RHS}|/\max(\|\mathrm{LHS}\|,\|\mathrm{RHS}\|)$
is checked against threshold $10^{-4}$. Across the
four generic-$\Omega$ test instances in
\texttt{test\_face\_model\_ddybe\_engine.py::
TestFaceDDYBEG2}
(\texttt{test\_ddybe\_generic\_omega},
\texttt{test\_ddybe\_different\_spectral},
\texttt{test\_ddybe\_different\_eta},
\texttt{test\_ddybe\_strong\_coupling}), the second-$\Omega$
spot check
\texttt{test\_ddybe\_generic\_omega\_second\_sample}, and the
diagonal-$\Omega$ sanity check
\texttt{test\_ddybe\_diagonal\_omega} (tier T3, consistency
with regime~(i)), the residual passes the stated threshold.
The $10^{-4}$ ceiling is an engineering budget chosen
conservatively above the float-$64$ accumulated roundoff of a
cascade of $8\times 8$ matrix products ($\sim 10^{-14}$)
and the $N=8$ truncation error of the genus-$2$ theta series;
it is a numerical threshold rather than a theoretical error term of the
form $O(\hbar^2)$.
The falsification-ladder test
\texttt{test\_ddybe\_ladder\_plateau} records the residual at
$N\in\{8,12,16\}$; the manuscript prior assertion that the
residual drops to $10^{-6}$ at $N=12$ is now a regression gate,
not a standalone empirical remark. This is \emph{numerical
evidence} for \texttt{conj:g2-ddybe} at finitely many parameter
points, not a proof and not a uniform bound across $\HHH_2$.
\end{proof}

\begin{remark}[What the numerical-evidence proposition does and does not establish]
\label{rem:g2p-regime-ii-scope}
Proposition~\ref{prop:g2-generic-ddybe-finite-window-evidence}
establishes: the face-model DDYBE
\eqref{eq:ddybe} holds to one part in $10^4$ across
$(z,w,\boldsymbol\lambda,\eta)$ parameter instances
sampled at two distinct $\Omega$ values
$\Omega_0, \Omega_1 \in \HHH_2 \setminus
\{\text{diagonal locus}\}$:
four parameter tuples at $\Omega_0$ plus one at $\Omega_1$.
The sampling is pointwise at two $\Omega$ values;
it leaves the full three-complex-dimensional
$\Omega$-dependence of the claim and the
non-separating boundary stratum where $\Omega_{22}\to
i\infty$ with $\Omega_{12}$ fixed nonzero
(Remark~\ref{rem:g2-nonseparating-untested}) untested. It establishes
neither an exact proof of \eqref{eq:ddybe} at
arbitrary $\Omega$ nor the stronger coupling
condition~\eqref{eq:ddybe-coupling} at chain level. The
latter remains conjectural and forms the open content of
\texttt{conj:g2-ddybe}.
\end{remark}

\begin{remark}[Tolerance ladder: engineering budget, not theoretical error]
\label{rem:tolerance-ladder}
The numerical claims distributed across the face-model
engine organise into four tiers. Each tier bound is an
\emph{engineering budget} (roundoff floor plus theta-series
truncation at $N$ lattice terms), not a theoretical error
term in the DDYBE itself: no finite-$\hbar$ bound of the form
$\|\mathrm{LHS}-\mathrm{RHS}\| \le c\,\hbar^p$ is proved.
The truncation error admits the standard exponential bound
$\varepsilon(N) \lesssim e^{-\pi N^2\,\min(\mathrm{Im}(\Omega))/2}$.
We state the tiers explicitly to eliminate source-of-truth
drift.
\begin{center}
\renewcommand{\arraystretch}{1.2}
\begin{tabular}{@{}llll@{}}
 \toprule
 Tier & Object & Tolerance & Justification \\
 \midrule
 (T1) & $\theta_1$ antisymmetry,
 weight-zero matrix entries
 & $<10^{-12}$ & exact algebraic identities in numpy \\
 (T2) & Genus-$1$ Felder DYBE
 & $<10^{-10}$ & $\theta$-series $N=30$ truncation \\
 (T3) & Diagonal-$\Omega$ degeneration to $g=1$
 & $<10^{-6}$ & $N=8$ genus-$2$ theta + algebraic cancellation \\
 (T4) & Generic-$\Omega$ DDYBE ($4$ instances)
 & $<10^{-4}$ & $N=8$ genus-$2$ theta, no cancellation \\
 \bottomrule
\end{tabular}
\end{center}
Tier~(T4) is the genuinely $g=2$ content; tiers~(T1)--(T3)
are preparatory or degeneration checks. The current face-model
verification suite contains $32$ tests across \texttt{TestThetaFunctions},
\texttt{TestBoltzmannWeightsG1}, \texttt{TestFaceDYBEG1},
\texttt{TestClassicalLimit}, \texttt{TestGenus2Theta},
\texttt{TestBoltzmannWeightsG2}, \texttt{TestFaceDDYBEG2}
(now including \texttt{test\_ddybe\_generic\_omega\_second\_sample}
and \texttt{test\_ddybe\_ladder\_plateau}),
\texttt{TestFaceSzegoThreeTerm} (new,
\texttt{test\_szego\_three\_term\_g2}),
and \texttt{TestFullSuite}).
\end{remark}

\begin{remark}[Numerical stability test]
\label{rem:tolerance-ladder-falsification}
The $N=8$ truncation of the $32$ face-model tests above is a
single finite-window snapshot. The regression test
\texttt{test\_ddybe\_ladder\_plateau} runs
\texttt{verify\_face\_ddybe\_g2} at $N\in\{8,12,16\}$ for two
distinct generic-$\Omega$ parameter instances and asserts that
the residual at $N\geq 12$ is below $10^{-6}$ and decreases
(or plateaus) from $N=8$ to $N=12$ to $N=16$. Any failure of the
$N\geq 12$ bound would indicate a finite-$\hbar$ DDYBE residual
that the $N=8$ bound $10^{-4}$ does not detect. Thus regime~(ii)
is finite-window numerical evidence at $N=8$ until the
$N\in\{12,16\}$ stability test is supplied; after a failed stability
test it becomes numerical evidence with explicit residual
$r_{\infty}>0$.
\end{remark}

\begin{remark}[Two distinct identities both due to Fay]
\label{rem:fay-versus-boltzmann-unitarity}
The engine function \texttt{verify\_unitarity\_g2} tests the
scalar product identity
\[
 \alpha\beta - \gamma\delta \;=\;
 \Theta_\delta((z-\eta)e_1\mid\Omega)\,/\,
 \Theta_\delta((z+\eta)e_1\mid\Omega),
\]
which is the Riemann-addition corollary governing Boltzmann-
weight unitarity. It is distinct from the normalised
Szeg\H{o}/prime-form Fay identity~\eqref{eq:fay-three-term}
(Fay 1973 Cor.~2.5), whose torus form contains the
$\vartheta_1'(0,\tau)^2$ prefactor and which governs the
bar-complex Stokes analysis in Remark~\ref{rem:why-a-not-b}.
Both are theta identities due to Fay; only~\eqref{eq:fay-three-term}
is the trisecant form. The bar-complex argument uses the
trisecant form in the bar-differential $d_{\mathrm{bar}}^2 = 0$
proof and the Boltzmann-unitarity form in the face-model
numerical checks.
\end{remark}

\begin{remark}[Non-separating degeneration: genuine open content]
\label{rem:g2-nonseparating-untested}
The non-separating degeneration
$\Omega_{22}\to i\infty$, $\Omega_{11}$ fixed, $\Omega_{12}$
\emph{fixed nonzero}, carries genuinely $g=2$ content: the
off-diagonal period $\Omega_{12}$ survives the limit and
encodes twisted modular data that is invisible to the
separating (diagonal-$\Omega$) degeneration. The current
engine suite does not test this boundary stratum. The
predicted limit is a genus-$1$ \emph{twisted} Felder DYBE
with Stokes data encoding $\Omega_{12}$; this prediction is
the genuine residual content of
\texttt{conj:g2-ddybe} and is a frontier target.
\end{remark}

% ========================================================
\section{Euler characteristic at degree~$2$}
\label{sec:g2-chi-minus-12}

\begin{corollary}[$\chi=-12$ from rank-$4$ KZB local system;
\ClaimStatusProvedHere]
\label{cor:g2-chi-minus-12}
\index{Euler characteristic!genus-$2$ degree-$2$ chiral homology}
On the once-punctured surface $\Sigma_2\setminus\{0\}$, the
KZB local system of rank $r=d^2=4$ (fibre $V\otimes V$ with
$V=\CC^2$) has twisted Euler characteristic
\begin{equation}\label{eq:g2p-euler}
 \chi\bigl(\Sigma_2\setminus\{0\},\cL^{(2)}_{\mathrm{KZB}}\bigr)
 \;=\; r\cdot\chi_{\mathrm{top}}(\Sigma_2\setminus\{0\})
 \;=\; 4\cdot(2-2g-s)\big|_{g=2,s=1}
 \;=\; 4\cdot(-3) \;=\; -12.
\end{equation}
For generic $\hbar$, $H^0=H^2=0$ and
$\dim H^1_{\mathrm{dR}}=12$ decomposes along isospin as
$H^1=\mathrm{Sym}^2 V\otimes\CC^3\;\oplus\;\bigwedge^2 V\otimes\CC^3$
(with $\dim\mathrm{Sym}^2 V=3$ giving $3\cdot 3=9$ and
$\dim\bigwedge^2 V=1$ giving $1\cdot 3=3$; total $12$).
\end{corollary}

\begin{proof}
Direct specialisation of
Proposition~\ref{prop:genus-g-euler-general} of
\texttt{higher\_genus\_modular\_koszul.tex} (which proves
$\chi=d^2(1-2g)$ on $\Sigma_g\setminus\{0\}$) to $d=2, g=2$.
\end{proof}

\begin{remark}[$\chi=-12$ is a topological invariant, not a
doubly-dynamical invariant]
\label{rem:chi-not-dd}
The value $\chi=-12$ is fully determined by
$(d,g,s)=(2,2,1)$ and does not depend on the dynamical
parameters $(\boldsymbol\lambda,\Omega_{12})$. The genuine
doubly-dynamical content of \texttt{conj:g2-ddybe} resides in
$R(\boldsymbol\lambda,\Omega,z)\in\End(V\otimes V)$ and
is invisible to the Euler-characteristic degree count. The
independent verification of $\chi=-12$ via Verlinde
($CB_{2,2}(1)=4$ at level $k=1$, matching the integrable
truncation of the $12$-dimensional generic $H^1$) is a
cross-check of the degree structure, not of the DDYBE.
\end{remark}

\begin{remark}[Integrable truncation]
\label{rem:integrable-truncation}
At integer level $k=1$, the degree-$2$ conformal block
count is $CB_{2,2}(1)=4$, from
Proposition~\ref{prop:g2-conformal-block-degree}. The
generic $H^1$ has dimension $12$; the integrable subspace
of dimension $4$ is the kernel of the $k+1$-fold $e_\theta$
singular vector, restricted to the physical sector.
The ratio $4/12 = 1/3$ is a consequence of the singular
vector count at level $1$ in the $\fsl_2$ minimal model,
not a doubly-dynamical invariant.
\end{remark}

% ========================================================
\section{Scope reconciliation}
\label{sec:g2-corrigenda}

\begin{remark}[Scope of the genus-$2$ DDYBE checks]
\label{rem:g2-corrigenda}
The genus-$2$ DDYBE checks separate into three statements:
\begin{enumerate}[label=\textup{(\alph*)}]
\item The diagonal/separating face-model degeneration is exact:
 Theorem~\ref{thm:g2-face-model-bypass-scope-restricted}
 reduces it to Felder's genus-$1$ DYBE.
\item The single-entry weight-zero check has agreement at
 $10^{-12}$; the generic-$\Omega$ DDYBE computation has numerical
 evidence at \texttt{relative < 1e-4} across five compute tests,
 recorded in
 Proposition~\ref{prop:g2-generic-ddybe-finite-window-evidence}.
\item The full generic-$\Omega$ DDYBE face-model statement remains
 conjectural:
 \texttt{conj:g2-ddybe} carries \texttt{\ClaimStatusConjectured}.
\item Fay~1973 Corollary~2.5 is the normalised
 Szeg\H{o}/prime-form trisecant identity and holds uniformly at
 all $g\ge 0$ by direct application. Its torus form contains the
 load-bearing $\vartheta_1'(0,\tau)^2$ prefactor. It is distinct
 from the four-theta-nulls quartic identity on the theta divisor
 and from the numerically tested Boltzmann-weight unitarity
 identity
 $\alpha\beta-\gamma\delta = \Theta(z-\eta)/\Theta(z+\eta)$.
\item There are $30$ face-model tests: four genuinely
 generic-$\Omega$ DDYBE tests at \texttt{relative < 1e-4}, one
 diagonal-$\Omega$ tier-(T3) consistency check against the genus-$1$
 Felder theorem, and $25$ tests for genus-$1$ DYBE, $\theta$
 antisymmetry, unitarity, crossing, classical limit, and degeneration.
\item The topological value $\chi=-12$ is not itself
 doubly-dynamical; by Remark~\ref{rem:chi-not-dd}, the
 doubly-dynamical content lives in
 $R(\boldsymbol\lambda,\Omega,z)$ and is currently the content of
 \texttt{conj:g2-ddybe}.
\end{enumerate}
\end{remark}

\begin{remark}[Engine label anchors]
\label{rem:engine-docstring-corrigenda}
The docstring of
\texttt{compute/lib/genus2\_ddybe\_engine.py} references
\texttt{prop:g2-nonsep-degen} and
\texttt{prop:g2-sep-degen}; both labels are now declared in
\texttt{higher\_genus\_modular\_koszul.tex}. This local DDYBE
chapter uses the following anchors for the same content:
\begin{itemize}
\item \texttt{prop:g2-nonsep-degen}: at the non-separating
 boundary $\Omega_{22}\to i\infty,\Omega_{12}\to 0$,
 the face-model $R$-matrix and DDYBE degenerate to a
 single genus-$1$ Felder DYBE in variable
 $\tau=\Omega_{11}$. This content is present in
 Remark~\ref{rem:g2p-g1-limit} above. It is not the
 off-diagonal non-separating frontier with $\Omega_{12}$ fixed
 nonzero, which remains open in
 Remark~\ref{rem:g2-nonseparating-untested}.
\item \texttt{prop:g2-sep-degen}: at the separating boundary
 $\Omega_{12}\to 0$ with $\Omega_{11},\Omega_{22}$ fixed,
 the DDYBE reduces to a single genus-$1$ Felder DYBE in
 variable $\tau=\Omega_{11}$ (the $\tau_2$-sector entering
 only as a cancelling $\theta_3(0|\tau_2)$ normalisation).
 This content is present in
 Theorem~\ref{thm:g2-face-model-bypass-scope-restricted}(i)
 and carries no genuinely genus-$2$ content.
\end{itemize}
These anchors state exactly the genus-$1$ Felder reductions used by
the degeneration computations, while the finite-window generic
off-diagonal tests are kept on the numerical-evidence surface
Proposition~\ref{prop:g2-generic-ddybe-finite-window-evidence}.
\end{remark}

% ========================================================
\section*{Verification anchors}

\begin{center}
\renewcommand{\arraystretch}{1.15}
\begin{tabular}{@{}lll@{}}
\toprule
Claim & Verification & Location \\
\midrule
Flat KZB connection on
 $\overline{\cM}_{2,n}\times\HHH_2$
 & CEE09 &
 \cref{thm:genus-2-kzb-connection-platonic} \\
Fay trisecant (normalised Szeg\H{o}/prime-form) at all $g$
 & Fay 1973 &
 \cref{thm:fay-trisecant-genus-2-specific} \\
DDYBE at diagonal $\Omega$ (exact)
 & Exact diagonal degeneration &
 \cref{thm:g2-face-model-bypass-scope-restricted}(i) \\
DDYBE at generic $\Omega$ (numerical, $<10^{-4}$, 4 tests)
 & Four generic-\(\Omega\) residual tests at \(<10^{-4}\) &
 \cref{prop:g2-generic-ddybe-finite-window-evidence} \\
DDYBE at separating boundary (factorised limit)
 & Exact factorised limit &
 \cref{thm:g2-face-model-bypass-scope-restricted}(ii) \\
$\chi=-12$ at degree~$2$, $g=2$, $s=1$
 & Euler-characteristic computation &
 \cref{cor:g2-chi-minus-12} \\
Full DDYBE~\eqref{eq:ddybe}--\eqref{eq:ddybe-coupling}
 at arbitrary $\Omega$
 & Open conjectural identity &
 \texttt{conj:g2-ddybe} \\
\bottomrule
\end{tabular}
\end{center}
